% =============================================================================
% §2 RELATED WORK & THE CITATION GAP
% Source of every paragraph below: docs/related-work-positioning.md (Week-0 gate,
% completed 16/07/2026). Those paragraphs were written as drop-in positioning
% text against VERIFIED metadata; do not paraphrase them into stronger claims.
%
% Related work = POSITIONING, not a list (playbook §6.5): lead with the
% integrated contribution, name and DISTINGUISH the near-rivals.
%
% ⚠️ FRAMING DECISION #3 (ke-hoach §4): the BIBLIOMETRIC table (N papers that do
% not cite an exact solution / do not measure Dice) goes to SUPPLEMENTARY. If
% the first table an editor sees is "we counted papers that failed to cite X",
% the paper is classified as meta-science and desk-rejected as out of scope.
% Table I of the manuscript is DP vs. exhaustive search (§3) — mathematics, not
% an accusation.
% =============================================================================
\section{Related Work and the Citation Gap}
\label{sec:related}

\subsection{Metaheuristic and Quantum-Inspired Multilevel Thresholding}

\TODO{2--3 paragraphs surveying the line of work, at the level of PRACTICE, not
of named authors. Cite from refs.bib only. Any count of papers ("N of M studies
report no Dice") comes from the coded bibliometric survey and may be stated only
once that survey exists with its two-coder Cohen's $\kappa$ --- see the
supplementary table and docs/RESULTS.md.}

% Recent, large, medical-image optimizer benchmark: the strongest checkable
% evidence for the thesis. Numbers in this paper's records are [PARTIAL] —
% re-open the full text before quoting any of them (positioning §7).
Al-Najdawi \emph{et al.}~\cite{alnajdawi_scirep} provide a large-scale, recent
benchmark of optimization algorithms coupled with Otsu's method on medical
images, reporting strong overlap scores while acknowledging in their limitations
the absence of ground-truth annotations. We treat this not as a competitor but
as evidence for the question we ask: the field continues to rank optimizers on
medical images by surrogate criteria. \TODO{quote their exact figures only after
re-opening the full text --- currently [PARTIAL], gated.}

% Positioning of the quantum-inspired branch.
% RED LINE (CLAUDE.md §1, lằn ranh 3): the words "quantum advantage", "quantum
% speedup", "quantum parallelism" are BANNED anywhere in this manuscript.
% Correct positioning: "an EDA-style probabilistic update rule", anchored to
% Platel et al., who showed QIEA is an EDA. Submission gate greps for this.
Quantum-inspired evolutionary algorithms are classical algorithms: their qubit
representation and rotation-gate update have been characterized as an
estimation-of-distribution update rule~\cite{platel_tevc}. We therefore describe
the quantum-inspired component throughout as an EDA-style probabilistic update
rule. No benefit arising from quantum computation is claimed or implied: such a
claim would require quantum hardware, which is not involved here, and the
algorithms studied are classical throughout.
% Phrased to avoid the literal banned string so the submission grep stays clean.

\subsection{The Exact Solution Already Exists}

% Drop-in paragraph, positioning §5. Cited in the ABSTRACT as well.
Exact, polynomial-time computation of optimal multilevel thresholds --- for
Otsu, Kapur, and Kittler--Illingworth criteria alike --- was given by Menotti
\emph{et al.}~\cite{menotti_ciarp}, with an $O((K-1)L^2)$ dynamic program.
Related exact formulations appear in Luessi \emph{et al.}~\cite{luessi_jei} and
Merzban and Elbayoumi~\cite{merzban_eswa}, the latter demonstrating that the
exact solver dominates metaheuristics for the Otsu criterion on natural images.
We claim \emph{no novelty} for the exact solver; it is the reference optimum
against which we benchmark. Nor do we claim novelty for benchmarking
metaheuristics against a global optimum: Hammouche \emph{et
al.}~\cite{hammouche_eaai} compared six metaheuristics against an exhaustive
% ⚠️ [PARTIAL] — the 1e-9 convergence criterion and their Tables 8--10 are
% recorded in docs/related-work-positioning.md §6 from an earlier full-text read;
% the metadata is [VERIFIED] but this number was NOT re-read from the publisher
% record. RE-OPEN THE FULL TEXT before submission, or drop the number and keep
% the qualitative statement.
global optimum, defining convergence as $|F - F_{\mathrm{opt}}| \le 10^{-9}$, a
stricter criterion than ours, sixteen years earlier. Two things distinguish our
study. First, those studies operate on gray-level images without clinical ground
truth, so a residual fitness gap cannot be projected onto a Dice-scored mask; we
make exactly that projection the primary endpoint. Second, Hammouche \emph{et
al.} already observe that some methods fail to reach the optimum for larger
threshold counts --- a finding we corroborate and reframe rather than contradict.

A dynamic-programming treatment that also \emph{selects} the number of
thresholds, including on medical images, has recently been given by Hegazy and
Gabr~\cite{hegazy_dp}. We claim novelty for neither the exact solver nor
data-driven selection of the threshold count; our use of dynamic programming is
purely as a reference optimum, and our treatment of the threshold count is to
measure what selecting it under surrogate criteria costs in clinical terms.

\subsection{Evaluation Standards for Segmentation}

Mousavirad \emph{et al.}~\cite{mousavirad_kbs} pose precisely our optimization
question --- whether population-based metaheuristics are effective for
multilevel thresholding --- but in a setting that structurally cannot detect the
phenomenon we study: the objective is Otsu between-class variance on natural
benchmark images without segmentation ground truth, so performance is measured
\emph{in the objective's own units}. They report that recent algorithms differ
substantially, some failing to reach the optimum.
% RECONCILIATION #1 (positioning §HOA GIAI #1) — this apparent contradiction
% with our equivalence prediction MUST be addressed here, not left for a
% reviewer to find. It is resolved by the axis of measurement, and it is turned
% into a testable result by the decoupling test (prereg §6/A4b, P2c).
Our contribution is orthogonal on two axes: we evaluate against an exact global
optimum rather than against each other, and we project every solution onto a
clinical binary mask with ground truth. A fitness gap that is real in variance
units need not survive that projection; whether it does is itself one of our
pre-specified tests, rather than an assumption.

Hegazy and Gabr~\cite{hegazy_metricbias} demonstrate, on natural images, that
PSNR and SSIM correlate systematically more with Otsu's objective than with
Kapur's, cautioning that reported gains may reflect this bias rather than
genuine improvements. Our corresponding analysis is a \emph{domain extension},
not an independent discovery: we move the argument to clinical ground truth on
brain MRI and add the threshold-count axis.

% RECONCILIATION #2 (positioning §HOA GIAI #2) — the ceiling belongs to
% Francois & Tinarrage. Stating this plainly, in Related Work, is the defense.
The claim that intensity thresholding has a performance ceiling on brain MRI is
not ours to make: François and Tinarrage~\cite{francois_jmiv} already report, on
BraTS FLAIR, an oracle best Dice of $0.83 \pm 0.18$, describing this explicitly
as an upper bound on the achievable performance within the considered
thresholding framework. We take this as an established result and make no claim
to establishing a ceiling; our contribution is its \emph{decomposition}
(§\ref{sec:ceiling}).
% The one quoted number in this manuscript that is NOT ours: it is a quotation
% from a [VERIFIED] full text (positioning §2), attributed in the same sentence.
% Everything else numeric must come from results/ via build_tables.py.

\subsection{The Gap}

\TODO{closing paragraph naming the gap in one sentence, written after §3 and
§8. It must be a statement about what has not been MEASURED, not a statement
about what other authors failed to do.}
